\subsection{Held-out generalization}\label{subsec:eval-heldout-generalization}

The held-out generalization test answers RQ5: whether the cascade configuration generalizes to papers it was never tuned on.It is run on a randomly selected subset of 15 papers, disjoint from the calibration cluster. The findings are scored with the identical matcher and matching threshold used on the calibration cluster (Subsection~\ref{subsec:eval-knowledge-extraction}): no tuning, threshold adjustment, or model selection is done on this held-out set. Therefore, the only intended difference is the set of papers. The held-out strict-F1 is reported in Chapter~\ref{chapter:Results}, alongside the in-sample calibration strict-F1; their difference is the generalization gap.

The test is evaluated against Opus-generated silver labels rather than expert annotations. It therefore measures whether the performance of the pipeline transfers to the held-out set relative to this silver standard, not whether the silver standard itself is clinically correct or complete. This limitation is mentioned in Subsection~\ref{subsec:eval-limitations}.